% ==============================================================================
% 公共题库: teaching/pool/解三角形/基础/q_triangle_dof_homogeneous.tex
% 难度: ★ 基础
% 知识点: 自由度分析(DOF=2, 形状锁定)、齐次条件、正弦定理边角互化
% ==============================================================================

\newcommand{\qTriangleDofHomogeneousStem}{%
在 \(\triangle ABC\) 中，内角 \(A, B, C\) 的对边分别为 \(a, b, c\)。已知 \(b+c=2a\)，且 \(5c\sin B = 6a\sin C\)。
\begin{enumerate}[(1)]
  \item 求三边长之比 \(a:b:c\)；
  \item 求 \(\cos B\) 的值。
\end{enumerate}%
}

\newcommand{\qTriangleDofHomogeneousAnswer}{%
(1) \(a:b:c = 5:6:4\)；\qquad (2) \(\cos B = \frac{1}{8}\)。
}

\newcommand{\qTriangleDofHomogeneousSolution}{%
由正弦定理角化边得 \(5cb = 6ac \implies b = \frac{6}{5}a\)。
代入 \(b+c=2a\) 得 \(c = \frac{4}{5}a\)。
故 \(a:b:c = 5:6:4\)。
由余弦定理 \(\cos B = \frac{a^2+c^2-b^2}{2ac} = \frac{1 + \frac{16}{25} - \frac{36}{25}}{\frac{8}{5}} = \frac{1}{8}\)。
}

\newcommand{\qTriangleDofHomogeneousDiagnosis}{%
【错因】未识别齐次条件，尝试盲目设具体边长，或无法完成边角比例翻译。
}
